\documentclass[11pt,a4paper]{article}

% ── Codificación y fuentes ────────────────────────────────────────────────────
\usepackage[utf8]{inputenc}
\usepackage[T1]{fontenc}
\usepackage{lmodern}
\usepackage[spanish]{babel}
\usepackage{amsmath}

% ── Geometría y espaciado ─────────────────────────────────────────────────────
\usepackage[top=2.5cm, bottom=2.5cm, left=2.8cm, right=2.8cm]{geometry}
\usepackage{setspace}
\setstretch{1.15}
\usepackage{parskip}

% ── Tablas ────────────────────────────────────────────────────────────────────
\usepackage{booktabs}
\usepackage{tabularx}
\usepackage{array}
\usepackage{longtable}
\usepackage{enumitem}

% ── Colores y cajas ───────────────────────────────────────────────────────────
\usepackage[dvipsnames]{xcolor}
\usepackage{tcolorbox}
\tcbuselibrary{skins, breakable}

% ── Cabeceras y pies de página ────────────────────────────────────────────────
\usepackage{fancyhdr}
\usepackage{lastpage}
\pagestyle{fancy}
\fancyhf{}
\fancyhead[L]{\small\textbf{Informe de Evaluación} --- Electrónica}
\fancyhead[R]{\small Tirzo Mata}
\fancyfoot[C]{\small Página \thepage\ de \pageref{LastPage}}
\renewcommand{\headrulewidth}{0.4pt}

% ── Hiperenlaces ──────────────────────────────────────────────────────────────
\usepackage[colorlinks=true, linkcolor=NavyBlue, urlcolor=NavyBlue]{hyperref}

% ── Paleta de colores personalizados ─────────────────────────────────────────
\definecolor{desempeno_excelente}{HTML}{1A6B2F}
\definecolor{desempeno_bueno}{HTML}{2E5A9C}
\definecolor{desempeno_suficiente}{HTML}{8B6914}
\definecolor{desempeno_deficiente}{HTML}{C0392B}
\definecolor{desempeno_insuficiente}{HTML}{7B241C}
\definecolor{grisFondo}{HTML}{F5F5F5}
\definecolor{azulTitulo}{HTML}{1B2A6B}
\definecolor{verdeFortaleza}{HTML}{1A6B2F}
\definecolor{naranjaMejora}{HTML}{A04000}

% ── Estilos de cajas ─────────────────────────────────────────────────────────
\tcbset{
  cajaTitulo/.style={
    enhanced, breakable,
    colback=azulTitulo!8, colframe=azulTitulo,
    fonttitle=\bfseries\large, coltitle=white,
    attach boxed title to top left={yshift=-2mm, xshift=4mm},
    boxed title style={colback=azulTitulo},
    top=6pt, bottom=6pt, left=8pt, right=8pt,
  },
  cajaFortaleza/.style={
    enhanced, breakable,
    colback=verdeFortaleza!6, colframe=verdeFortaleza!60,
    fonttitle=\bfseries, coltitle=verdeFortaleza,
    top=4pt, bottom=4pt, left=8pt, right=8pt,
  },
  cajaMejora/.style={
    enhanced, breakable,
    colback=naranjaMejora!6, colframe=naranjaMejora!60,
    fonttitle=\bfseries, coltitle=naranjaMejora,
    top=4pt, bottom=4pt, left=8pt, right=8pt,
  },
  cajaRecomendacion/.style={
    enhanced, breakable,
    colback=grisFondo, colframe=gray!50,
    fonttitle=\bfseries, coltitle=black,
    top=4pt, bottom=4pt, left=8pt, right=8pt,
  },
}

% ─────────────────────────────────────────────────────────────────────────────
\begin{document}

% ══ PORTADA ══════════════════════════════════════════════════════════════════
\begin{center}
  {\Large\bfseries\color{azulTitulo} ELECTRÓNICA --- Evaluación de Exámenes}\\[4pt]
  {\large Informe de Evaluación Preliminar}\\[12pt]
  \rule{\linewidth}{1.2pt}\\[6pt]
  {\LARGE\bfseries Tirzo Mata}\\[4pt]
  \rule{\linewidth}{0.6pt}
\end{center}

% ══ FICHA TÉCNICA ════════════════════════════════════════════════════════════
\vspace{4pt}
\begin{tcolorbox}[cajaTitulo, title=Datos del informe]
\begin{tabular}{@{}llll@{}}
  \textbf{Estudiante:}  & Tirzo Mata  &
  \textbf{Fecha:}       & 2026-06-27 \\[4pt]
  \textbf{Archivo PDF:} & \texttt{examen\_Tirzo\_Mata.pdf}  &
  \textbf{Páginas:}     & 4 \\
\end{tabular}
\end{tcolorbox}

% ══ RESULTADO GLOBAL ═════════════════════════════════════════════════════════
\vspace{6pt}
\begin{center}
  \begin{tcolorbox}[
    enhanced,
    colback=white, colframe=desempeno_excelente,
    width=0.62\linewidth,
    boxrule=2pt,
    halign=center, valign=center,
    top=8pt, bottom=8pt,
  ]
    \centering
    {\large\bfseries Puntaje sugerido}\\[6pt]
    {\Huge\bfseries\color{desempeno_excelente} 10.0/10}\\[6pt]
    {\large Nivel de desempeño: \textbf{\color{desempeno_excelente} Excelente}}
  \end{tcolorbox}
\end{center}

% ══ RESUMEN GENERAL ══════════════════════════════════════════════════════════
\section*{Resumen general}
El estudiante demuestra un dominio sobresaliente de los conceptos fundamentales de la teoría de circuitos. Resolvió de manera impecable los tres problemas planteados: simplificación de resistencia equivalente, análisis de mallas y análisis nodal. El desarrollo procedimental es sumamente ordenado, claro y matemáticamente exacto.

% ══ FORTALEZAS Y ÁREAS DE MEJORA ════════════════════════════════════════════
\vspace{4pt}
\begin{minipage}[t]{0.48\linewidth}
  \begin{tcolorbox}[cajaFortaleza, title={Fortalezas}]
\begin{itemize}[leftmargin=1.5em, itemsep=2pt]
  \item Excelente nivel de orden y claridad en el desarrollo de los diagramas y las ecuaciones.
  \item Uso correcto y consistente de las unidades de medida (Ohmios, Amperios, Voltios) en los resultados finales.
  \item Dominio impecable del álgebra para la resolución de sistemas de ecuaciones lineales.
  \item Comprensión profunda de las leyes de Kirchhoff (LVK y LCK) aplicadas a mallas y nodos.
\end{itemize}
  \end{tcolorbox}
\end{minipage}
\hfill
\begin{minipage}[t]{0.48\linewidth}
  \begin{tcolorbox}[cajaMejora, title={Areas de mejora}]
\begin{itemize}[leftmargin=1.5em, itemsep=2pt]
  \item Como detalle menor de presentación, se podría sugerir evitar tachaduras o encimados leves en los cálculos de la página 3, aunque no afectan en absoluto la legibilidad ni la corrección del ejercicio.
\end{itemize}
  \end{tcolorbox}
\end{minipage}

% ══ OBSERVACIONES POR PREGUNTA ═══════════════════════════════════════════════
\section*{Observaciones por pregunta}

\begin{longtable}{>{\bfseries}p{2.8cm} p{1.6cm} p{7.0cm} p{2.8cm}}
  \toprule
  \textbf{Pregunta} & \textbf{Puntaje} & \textbf{Evaluación} & \textbf{Errores detectados} \\
  \midrule
  \endhead
  \bottomrule
  \endfoot
    \textbf{Pregunta 1: Resistencia Equivalente} & 3.3/3.3 & El estudiante simplifica el circuito paso a paso de manera excelente. Identifica correctamente las asociaciones en paralelo ($R_2 \parallel R_3$ y luego $R_{eq2} \parallel R_4$) y en serie ($R_1 + R_{eq1}$ y finalmente $R_{eq3} + R_5$). Dibuja los circuitos equivalentes intermedios, lo cual facilita enormemente el seguimiento del procedimiento. & \textit{---} \\
    \midrule
    \textbf{Pregunta 3: Análisis de Mallas} & 3.4/3.4 & Plantea correctamente las ecuaciones de malla utilizando la Ley de Tensiones de Kirchhoff (LVK). Resuelve el sistema de ecuaciones de $2 \times 2$ mediante el método de eliminación de forma impecable, obteniendo los valores correctos para las corrientes de malla $I_1 \approx 1.42\text{ A}$ e $I_2 \approx 0.16\text{ A}$. Finalmente, calcula con precisión la corriente que circula por el resistor central de $5\,\Omega$ ($I_R \approx 1.26\text{ A}$). & \textit{---} \\
    \midrule
    \textbf{Pregunta 2: Análisis Nodal} & 3.3/3.3 & Plantea de forma correcta las ecuaciones de nodo utilizando la Ley de Corrientes de Kirchhoff (LCK) para los nodos A y B. Realiza la sustitución algebraica de manera directa y elegante, obteniendo los voltajes de nodo correctos: $V_B = 5\text{ V}$ y $V_A = 7.5\text{ V}$. & \textit{---} \\
    \midrule
\end{longtable}

% ══ ERRORES DETECTADOS ═══════════════════════════════════════════════════════
\section*{Análisis de errores}

\subsection*{Errores conceptuales}
\textit{Ninguno identificado.}

\subsection*{Errores procedimentales}
\textit{Ninguno identificado.}

% ══ RECOMENDACIONES AL ESTUDIANTE ════════════════════════════════════════════
\section*{Retroalimentación para el estudiante}
\begin{tcolorbox}[cajaRecomendacion, title={Dirigido a: Tirzo Mata}]
¡Excelente trabajo! Tu examen muestra un nivel de preparación y orden sobresaliente. Continúa con esa misma rigurosidad metodológica y claridad en el desarrollo de tus circuitos en los siguientes temas de la materia.
\end{tcolorbox}

% ══ NOTA PARA EL PROFESOR ════════════════════════════════════════════════════
\section*{Nota para el profesor}
\begin{tcolorbox}[
  enhanced, breakable,
  colback=yellow!8, colframe=yellow!60!black,
  fonttitle=\bfseries, coltitle=black,
  title={Observaciones para revisión manual},
  top=4pt, bottom=4pt, left=8pt, right=8pt,
]
El examen está perfecto. No presenta errores conceptuales ni de cálculo. Los diagramas de simplificación de la pregunta 1 y el desarrollo algebraico de los sistemas de ecuaciones en las preguntas 2 y 3 están completamente detallados, lo que demuestra un trabajo honesto y de alta calidad.
\end{tcolorbox}

% ══ PIE DE INFORME ════════════════════════════════════════════════════════════
\vspace{12pt}
\begin{center}
  \footnotesize\color{gray}
  Informe generado automáticamente por el sistema de evaluación con IA (gemini-3.5-flash) y soporte LaTex. \\
  Este documento es una evaluación \textbf{preliminar} para ser revisado por el profesor antes de ser entregado al 
  estudiante. \\
  Generado el 2026-06-27 \textbullet\ ELECTRÓNICA --- Sistema de Evaluación Académica.
  \end{center}

\end{document}
